\documentclass[../../../include/open-logic-section]{subfiles}

\begin{document}

\olfileid{sth}{choice}{intro}

\olsection{سريزه}

په \crefrange{sth:cardinals::chap}{sth:card-arithmetic::chap} کښې مو
اصلي عددونه د ترتيبي عددونو په توګه په ګڼلو سره د اصلي عددونو
نظريه جوړه کړه۔ دا لاره د ښه ترتيب پر اصل ولاړه ده۔ څرګندېږي
چې ښه ترتيب له يوه بل اصل، يعنې د انتخاب له اصل سره، همارز
دے؛ د دې اصل د منلو په اړه جدي فلسفي بحث هم شوے دے۔
په دې څپرکي کښې مو پوښتنې دا دي: دا اصل څنګه کارول کېږي،
او ايا توجيه کېدلے شي؟

\end{document}